% 分野別類題: 波動方程式の有限区間問題（両端固定・変数分離法）
% u_tt = c^2 u_xx （0 < x < L）を変数分離とフーリエ級数で解く 3 問。
% コンパイル: TEXINPUTS="../../../00_common:$TEXINPUTS" latexmk -xelatex problems.tex
\documentclass[a4paper,11pt]{article}
\usepackage{preamble}

\begin{document}

\kamokumei{類題演習 --- 波動方程式の有限区間問題（変数分離法）}

\shijibun{次の波動方程式の初期値・境界値問題を、変数分離法（フーリエ級数展開）により解け。導出過程を示すこと。}

\shomon{問1.}{長さ $L$ の弦の振動を表す次の初期値・境界値問題を解け。}
\[
\frac{\partial^{2} u}{\partial t^{2}} = c^{2}\frac{\partial^{2} u}{\partial x^{2}} \qquad (0 < x < L,\ t > 0)
\]
\[
u(0,t) = 0,\quad u(L,t) = 0 \qquad (t > 0)
\]
\[
u(x,0) = \sin\frac{\pi x}{L}, \qquad \frac{\partial u}{\partial t}(x,0) = 0 \qquad (0 \le x \le L)
\]

\shomon{問2.}{長さ $L$ の弦の振動を表す次の初期値・境界値問題を解け。}
\[
\frac{\partial^{2} u}{\partial t^{2}} = c^{2}\frac{\partial^{2} u}{\partial x^{2}} \qquad (0 < x < L,\ t > 0)
\]
\[
u(0,t) = 0,\quad u(L,t) = 0 \qquad (t > 0)
\]
\[
u(x,0) = f(x), \qquad \frac{\partial u}{\partial t}(x,0) = 0 \qquad (0 \le x \le L)
\]
ただし初期変位は
\[
f(x) =
\begin{cases}
\dfrac{2h}{L}\,x & \left(0 \le x \le \dfrac{L}{2}\right) \\[2mm]
\dfrac{2h}{L}\,(L-x) & \left(\dfrac{L}{2} \le x \le L\right)
\end{cases}
\]
（弦の中点を高さ $h$ だけ持ち上げて静かに放した状態）で与えられるとする。フーリエ正弦級数の係数 $b_{n}$ を求め、解 $u(x,t)$ を級数の形で表せ。

\shomon{問3.}{長さ $L$ の弦の振動を表す次の初期値・境界値問題を解け。}
\[
\frac{\partial^{2} u}{\partial t^{2}} = c^{2}\frac{\partial^{2} u}{\partial x^{2}} \qquad (0 < x < L,\ t > 0)
\]
\[
u(0,t) = 0,\quad u(L,t) = 0 \qquad (t > 0)
\]
\[
u(x,0) = 0, \qquad \frac{\partial u}{\partial t}(x,0) = v_{0}\sin\frac{3\pi x}{L} \qquad (0 \le x \le L)
\]
（初期変位は $0$、初期速度が $v_{0}\sin\dfrac{3\pi x}{L}$ で与えられる、いわゆる「打弦」型の初期条件。）解 $u(x,t)$ を求めよ。

\end{document}
